% Options for packages loaded elsewhere
\PassOptionsToPackage{unicode}{hyperref}
\PassOptionsToPackage{hyphens}{url}
\PassOptionsToPackage{dvipsnames,svgnames,x11names}{xcolor}
\documentclass[
  10pt,
  letterpaper,
]{article}
\usepackage{xcolor}
\usepackage[margin=1in]{geometry}
\usepackage{amsmath,amssymb}
\setcounter{secnumdepth}{5}
\usepackage{iftex}
\ifPDFTeX
  \usepackage[T1]{fontenc}
  \usepackage[utf8]{inputenc}
  \usepackage{textcomp} % provide euro and other symbols
\else % if luatex or xetex
  \usepackage{unicode-math} % this also loads fontspec
  \defaultfontfeatures{Scale=MatchLowercase}
  \defaultfontfeatures[\rmfamily]{Ligatures=TeX,Scale=1}
\fi
\usepackage{lmodern}
\ifPDFTeX\else
  % xetex/luatex font selection
\fi
% Use upquote if available, for straight quotes in verbatim environments
\IfFileExists{upquote.sty}{\usepackage{upquote}}{}
\IfFileExists{microtype.sty}{% use microtype if available
  \usepackage[]{microtype}
  \UseMicrotypeSet[protrusion]{basicmath} % disable protrusion for tt fonts
}{}
\makeatletter
\@ifundefined{KOMAClassName}{% if non-KOMA class
  \IfFileExists{parskip.sty}{%
    \usepackage{parskip}
  }{% else
    \setlength{\parindent}{0pt}
    \setlength{\parskip}{6pt plus 2pt minus 1pt}}
}{% if KOMA class
  \KOMAoptions{parskip=half}}
\makeatother
\usepackage{listings}
\newcommand{\passthrough}[1]{#1}
\definecolor{ocamlkeyword}{RGB}{0,80,150}
\definecolor{ocamlmodule}{RGB}{110,55,130}
\definecolor{ocamlcomment}{RGB}{80,125,75}
\definecolor{ocamlstring}{RGB}{160,70,50}
\definecolor{ocamlbackground}{RGB}{248,248,246}
\definecolor{ocamlframe}{RGB}{220,220,216}
\lstdefinelanguage{Caml}{
  morekeywords={
    and,as,assert,begin,class,constraint,do,done,downto,else,end,exception,
    external,false,for,fun,function,functor,if,in,include,inherit,initializer,
    lazy,let,let%op,let%cd,match,method,module,mutable,new,nonrec,object,of,
    open,or,private,rec,sig,struct,then,to,true,try,type,val,virtual,when,while,
    with
  },
  morekeywords=[2]{
    Assignments,Backends,Context,Float,Ir,NTDSL,Option,PDSL,Shape,TDSL,Tensor,
    Tnode
  },
  sensitive=true,
  morecomment=[s]{(*}{*)},
  morestring=[b]",
  alsoletter={\%,',_,!,?,.}
}
\lstset{defaultdialect=[5.3]Lua}
\lstset{defaultdialect=[x86masm]Assembler}
\lstset{
  basicstyle=\small\ttfamily,
  keywordstyle=\bfseries\color{ocamlkeyword},
  keywordstyle=[2]\color{ocamlmodule},
  commentstyle=\itshape\color{ocamlcomment},
  stringstyle=\color{ocamlstring},
  backgroundcolor=\color{ocamlbackground},
  breakatwhitespace=false,
  breaklines=true,
  columns=flexible,
  frame=single,
  framerule=0.3pt,
  rulecolor=\color{ocamlframe},
  keepspaces=true,
  numberstyle=\scriptsize\color{gray},
  showstringspaces=false,
  upquote=true,
  xleftmargin=0.6em,
  xrightmargin=0.6em,
  framesep=0.45em
}
\usepackage{longtable,booktabs,array}
\usepackage{calc} % for calculating minipage widths
% Correct order of tables after \paragraph or \subparagraph
\usepackage{etoolbox}
\makeatletter
\patchcmd\longtable{\par}{\if@noskipsec\mbox{}\fi\par}{}{}
\makeatother
% Allow footnotes in longtable head/foot
\IfFileExists{footnotehyper.sty}{\usepackage{footnotehyper}}{\usepackage{footnote}}
\makesavenoteenv{longtable}
\setlength{\emergencystretch}{3em} % prevent overfull lines
\providecommand{\tightlist}{%
  \setlength{\itemsep}{0pt}\setlength{\parskip}{0pt}}
\providecommand{\llbracket}{\mathopen{[\![}}
\providecommand{\rrbracket}{\mathclose{]\!]}}
\usepackage{bookmark}
\IfFileExists{xurl.sty}{\usepackage{xurl}}{} % add URL line breaks if available
\urlstyle{same}
\hypersetup{
  pdftitle={Neural Networks Without Shape Boilerplate: An OCaml DSL Case Study},
  colorlinks=true,
  linkcolor={blue},
  filecolor={Maroon},
  citecolor={Blue},
  urlcolor={blue},
  pdfcreator={LaTeX via pandoc}}

\title{Neural Networks Without Shape Boilerplate: An OCaml DSL Case
Study}
\author{}
\date{}

\begin{document}
\maketitle

\begin{abstract}

Modern neural network code is often cluttered with shape bookkeeping:
explicit calls to reshape, unsqueeze, expand, transpose, reductions, and
integer-axis arguments. OCANNL combines an embedded DSL that removes
this boilerplate, with an end-to-end compiler with GPU backends. Users
define tensor operations as OCaml functions using concise notation
generalizing the Einstein summation convention. Concatenation is an
expression over indices e.g.~\passthrough{\lstinline!i\^j!}, convolution
is a contraction with affine operand addressing
e.g.~\passthrough{\lstinline!stride*i + dilation*k!}. A sequence of axes
can be captured by a single variable,
e.g.~\passthrough{\lstinline!..batch..!} or
\passthrough{\lstinline!..activations..!}, sandwiched between leading
and trailing axes; with up to three such variables per shape. We call
these \emph{row variables}. The paper showcases three examples: tensor
expressions core multi-head attention, rank-polymorphic 2D convolution;
and tensor computation for Stochastic Gradient Descent with momentum and
Nesterov-inspired correction. OCANNL performs broadcast-aware global
bidirectional shape inference and derives loop-index maps for code
generation. We formalize the inference problem for the core calculus
(excluding affine and concatenation) and show properties of OCANNL's
solver (proofs in the appendix). OCANNL's compiler inlines computations
to avoid materializing intermediate tensors, and performs common
subexpression elimination. This paper does not showcase benchmarks nor
argue for performance, leaving that to follow-up work and the Q\&A.

\end{abstract}

\section{Shapes, Tensor Expressions, Syntax Extensions and
Examples}\label{shapes-tensor-expressions-syntax-extensions-and-examples}

OCANNL shapes have three kinds (sequences) of axes, from outermost:
batch, output, input. The kind designation is by convention (not
enforced, but facilitated). In specifications, the syntax resembles
programming language types:
\passthrough{\lstinline!batch | input -> output!}. For example: tensor
multiplication \passthrough{\lstinline!*!} reduces the input axes of the
operator tensor (matrix etc.) with output axes of the operand tensor
(vector or matrix etc.), while broadcasting or treating pointwise the
batch axes.

A \emph{tensor node} \passthrough{\lstinline!Tnode.t!} is an identity
that also carries dimensionality, precision and padding information, but
neither data (memory) nor computation. Computations themselves
(assignments, type \passthrough{\lstinline!Assignments.comp!}) are
expressed in terms of tensor nodes, which makes them hardware (backend)
agnostic.

\emph{Tensor expressions} \passthrough{\lstinline!Tensor.t!} are the
primary type of the OCANNL frontend. They combine an inferrable shape, a
value node and \emph{forward computation} to update it, and optionally a
gradient node and \emph{backprop computation} to propagate the gradient
to its component tensors. OCANNL tracks initialization via a field
\passthrough{\lstinline!Tensor.t.params!} that's a set of tensors.
OCANNL frontend does not have a separate abstraction for machine
learning models: any tensor can be considered to be a model, whose
trainable parameters are the differentiable subset of its
\passthrough{\lstinline!params!}.

OCANNL introduces two syntaxes. Extension point
\passthrough{\lstinline!\%op!} creates tensor expressions, or OCaml
functions returning tensor expressions which we call tensor
\emph{operations}. Extension point \passthrough{\lstinline!\%cd!}
creates tensor computations. Both extensions support inline definitions
of new tensors via OCaml's record syntax. For example:
\passthrough{\lstinline!\{ w1 = kaiming normal1 () \}!} inside
\passthrough{\lstinline!\%op!} introduces tensor
\passthrough{\lstinline!w1!} with initialization expression
\passthrough{\lstinline!kaiming normal1 ()!};
\passthrough{\lstinline!\{ sgd\_momentum \}!} inside
\passthrough{\lstinline!\%cd!} below introduces a (non-differentiable,
no initialization) tensor \passthrough{\lstinline!sgd\_momentum!};
\passthrough{\lstinline!\{ w\_q \}!} inside
\passthrough{\lstinline!\%op!} below introduces a tensor with default
initialization (e.g.~centered uniform distribution).

At the heart of OCANNL are indexing specifications for expressing tensor
computations. They generalize the Einstein summation convention: indices
missing from the result are reduced over. The specification syntax uses
\passthrough{\lstinline!;!} to separate arguments and
\passthrough{\lstinline!=>!} to separate out the result. Ellipsis
\passthrough{\lstinline!...!} in the specifications are expanded
contextually as either \passthrough{\lstinline!..batch..!},
\passthrough{\lstinline!..input..!} or
\passthrough{\lstinline!..output..!} row variables. Assignments in the
\passthrough{\lstinline!\%cd!} syntax specify the unary or binary
arithmetic and the accumulation arithmetic explicitly. For the
\passthrough{\lstinline!\%op!} syntax, we have mixfix operators
combining the arithmetic semantics, for example:
\passthrough{\lstinline!++!} is identity with additive accumulation,
\passthrough{\lstinline!+*!} is multiplication with additive
accumulation. These operators also support \emph{variable capture} for
both dimensions and rows -- the variables from a trailing string list
are introduced into scope earlier than they first appear (simiarly to
inline definitions). The captured variables can be used for both shape
constraints \passthrough{\lstinline!Shape.set\_dim!},
\passthrough{\lstinline!Shape.equal\_dims!} and converted to scalars via
\passthrough{\lstinline!dim!} (for row variables, the value is the
product of the dimensions of the axes).

\begin{lstlisting}[language=Caml]
let%op softmax ~spec ?(temperature = 1.0) () =
  let spec = spec ^ " => " ^ replace_alphanum_by_0 spec in
  fun x ->
    let x_scaled = if Float.(temperature <> 1.0) then x /. !.temperature else x in
    let max_vals = x_scaled @^^ spec in
    let exp_vals = exp (x_scaled - max_vals) in
    exp_vals /. (exp_vals ++ spec)

let%op multi_head_attention ~num_heads ~d_k ~d_v () x =
  let q = { w_q } * x in
  let k = { w_k } * x in
  let v = { w_v } * x in
  let scores =
    (q +* k " ... s | h d; ... t | h d => ... s | t -> h" [ "h"; "d" ])
    /. sqrt (dim d)
  in
  Shape.set_dim h num_heads;
  Shape.set_dim d d_k;
  Shape.set_dim e d_v;
  let attn_weights = softmax ~spec:" ... | t -> ..." () scores in
  { w_o } *
    (attn_weights +* v
       " ... s | t -> h; ... t | h e => ... s | h e" [ "e" ])
\end{lstlisting}

\begin{quote}
\textbf{Figure 1. Core multi-head attention in OCANNL.} Axes: query
position \passthrough{\lstinline!s!}, key/value position
\passthrough{\lstinline!t!}, head \passthrough{\lstinline!h!}, key width
\passthrough{\lstinline!d!}, and value width
\passthrough{\lstinline!e!}. Batch rank, parameter shapes, score shape,
contractions, and loop-index maps are inferred.
\passthrough{\lstinline!@\^!} is infix for
\passthrough{\lstinline!max!}, and \passthrough{\lstinline!@\^\^!} is
the \passthrough{\lstinline!max!} reducing equivalent of
\passthrough{\lstinline!++!}.
\end{quote}

\begin{lstlisting}[language=Caml]
let%op conv2d ~label ?(kernel_size = 3) ?(stride = 1)
    ?(use_padding = true) ?out_channels () x =
  Shape.set_dim kh kernel_size;
  Shape.set_dim kw kernel_size;
  Option.iter out_channels ~f:(Shape.set_dim oc);
  x +* { kernel }
       "... | stride*oh + kh, stride*ow + kw, ..ic..;
             kh, kw, ..ic.. -> ..oc..
        => ... | oh, ow, ..oc.."
       [ "kh"; "kw"; "oc" ]
  + { bias = 0. }
\end{lstlisting}

\begin{quote}
\textbf{Figure 2. Rank-polymorphic 2D convolution.} The input is
addressed at affine spatial positions
\passthrough{\lstinline!stride*oh+kh!} and
\passthrough{\lstinline!stride*ow+kw!}. The kernel axes
\passthrough{\lstinline!kh!} and \passthrough{\lstinline!kw!}, together
with the input-channel row \passthrough{\lstinline!..ic..!}, are
contracted. The context row \passthrough{\lstinline!...!} and
output-channel row \passthrough{\lstinline!..oc..!} are inferred and may
have arbitrary rank.
\end{quote}

\begin{lstlisting}[language=Caml]
let sgd_one ~learning_rate ?(momentum = 0.0) ?(weight_decay = 0.0)
    ?(nesterov = false) p =
  [%cd
    { sgd_delta } =: p.grad + (!.weight_decay *. p);
    if Float.(momentum > 0.0) then (
      { sgd_momentum } =:
        (!.momentum *. sgd_momentum) + sgd_delta;
      if nesterov then sgd_delta =+ !.momentum *. sgd_momentum
      else sgd_delta =: sgd_momentum);
    p =- learning_rate * sgd_delta ~logic:"."]
\end{lstlisting}

\begin{quote}
\textbf{Figure 3. Stochastic Gradient Descent with momentum.} Introduces
intermediate state-tracking tensors in a computation. Optional
Nesterov-inspired correction.
\end{quote}

OCANNL implements coproduct of axes \passthrough{\lstinline!\^!},
generalizing concatenation. It is unlocked by the n-ary operator
\passthrough{\lstinline!++\^!}. We extend the solution space for
dimensions with dimension-0, the unit of the coproduct; it is not a
valid shape dimension (actual axes never end up at dimension-0). The
operator \passthrough{\lstinline!++\^!} reduces additively (users can
define their own variants).

\begin{longtable}[]{@{}
  >{\raggedright\arraybackslash}p{(\linewidth - 4\tabcolsep) * \real{0.3333}}
  >{\raggedright\arraybackslash}p{(\linewidth - 4\tabcolsep) * \real{0.4000}}
  >{\raggedright\arraybackslash}p{(\linewidth - 4\tabcolsep) * \real{0.2667}}@{}}
\toprule\noalign{}
\begin{minipage}[b]{\linewidth}\raggedright
PyTorch/TensorFlow
\end{minipage} & \begin{minipage}[b]{\linewidth}\raggedright
OCANNL
\end{minipage} & \begin{minipage}[b]{\linewidth}\raggedright
Notes
\end{minipage} \\
\midrule\noalign{}
\endhead
\bottomrule\noalign{}
\endlastfoot
\passthrough{\lstinline!torch.cat([a, b], dim=0)!} &
\passthrough{\lstinline!(a, b) ++\^ "x, ...; y, ... => x\^y, ..."!} &
Concatenate tensors \\
\passthrough{\lstinline!torch.stack([a, b], dim=0)!} &
\passthrough{\lstinline![a; b]!} & Stack with new axis (block tensor
syntax) \\
\passthrough{\lstinline!x[:n]!} (prefix slice) &
\passthrough{\lstinline!x ++\^ "a\^b => a"!} & Extract prefix (size
inferred) \\
\passthrough{\lstinline!x[n:]!} (suffix slice) &
\passthrough{\lstinline!x ++\^ "a\^b => b"!} & Extract suffix (size
inferred) \\
\passthrough{\lstinline!x[:-3]!} (all but last 3) &
\passthrough{\lstinline!x ++\^ "a\^3 => a"!} & Drop last 3 elements \\
\end{longtable}

\section{Shapes and Projections: design choices and
claims}\label{shapes-and-projections-design-choices-and-claims}

We say ``row'' for analogy with type systems, but our shape rows are
sequences of axes rather than records. OCANNL doesn't have a
``full-blown'' parametric polymorphism because we don't generalize
inferred shapes into shape schemes, but it does have ``poor man's''
parametric polymophism because we generate fresh shape variables for
shape constraints derived at tensor operation call sites. Thus OCANNL
has two forms of rank polymorphism: structural (subtyping) via
broadcasting and parametric via row variables.

The motivation to have left flank axes in shape (row) terms first came
from the slice operation -- for example when subdividing data into
hierarchical minibatches (like in data parallel batched training). This
impacts broadcasting: rather assuming that implicit axes from
broadcasting land to the left of the ``left flank'', we generalize
broadcasting to happen in the middle between left and right flanks of
already-determined axes of a shape row. This allows uniform handling of
broadcast axes and inferred axes in inference (removing a major source
of non-determinism).

However, consistently applying the semantics that falls out of the
structural rank-polymorphism (subtyping polymorphism, formally
\(\sqsubseteq\) in the Appendix) to the indexing (aka. \emph{einsum})
specifications seen in examples above would be both counter-intuitive
and overly restrictive: the broadcasting behavior of the operands would
have to match the position of the row variables in the specification. To
relax this, we introduce a \emph{flat row equivalence} relation
(formally \(\approx\) in the Appendix), it preserves rank but erases the
broadcasting position. Except for relaxing broadcast position checking,
\(C \approx T^{\mathrm{lhs}}\) is stricter than
\(C \sqsubseteq T^{\mathrm{lhs}}\) would be (see Appendix section 3),
improving shape propagation through inference.

In the Appendix, we provide formal semantics for the core shape
constraints (for simplicity, excluding affine indexing and
concatenation), present the shape and projections inference algorithms
(where projections represent tensor indexing for tensor assignment
loops), and prove their properties. Proofs are available online. In
particular, TR Theorem 5.4, TR Theorem 5.9, and TR Proposition 6.2,
together state inference termination and soundness. Inference is
declarative: the result does not depend on the ordering of the
constraints. The shape inference problem is non-principal, and the
algorithm is in general incomplete. Should the incompleteness ever
manifest in practice, users can provide additional constraints via for
example \passthrough{\lstinline!Shape.infer\_equal!}, or variable
capture and \passthrough{\lstinline!Shape.set\_equal!}. Example of
incompleteness from TR Example 6.3, where our inference fails although
constraints are satisfiable, with \(\alpha,\beta\) coming from leaf
tensors and \(\gamma\) result tensor:

\[
\Phi_2 = \{3_b \sqsubseteq \alpha,\ 5_b \sqsubseteq \beta,\ \gamma \sqsubseteq \alpha,\ \gamma \sqsubseteq \beta\}
\]

Failing here is the more useful outcome, relative to guessing either
\(\alpha = 1_\emptyset, \gamma = 5_b\) or
\(\beta = 1_\emptyset, \gamma = 3_b\). Another example stems from the
leniency in the semantics of row equivalence. See Remark after
Proposition 2.9 in the appendix:

\[
\Phi = \{\,\langle\rho\rangle \approx [\,]\cdot\diamond\cdot[3,5],\ \ [3]\cdot\diamond\cdot[9,5] \sqsubseteq \langle\rho\rangle\,\}
\]

Solution \(\gamma\rho = [3]\cdot\diamond\cdot[5]\) would satisfy both
constraints (\([3,9,5]\) vs \([3,1_\emptyset,5]\)). OCANNL fails here by
comitting to \(\gamma\rho = [\,]\cdot\diamond\cdot[3,5]\).

\subsection{Dimension basis instead of axis
name}\label{dimension-basis-instead-of-axis-name}

OCANNL does not support attaching axis names to tensors, as both the
semantics and ergonomics of indexing are sequence based. Instead OCANNL
supports attaching bases (arbitrary identifiers) to dimensions.
Dimensions of different bases are incompatible even when they're of the
same size. \(1_b\) is not broadcastable for any basis \(b\) except for
\(1_\emptyset\) (basis \(\emptyset\) is the identifier
\passthrough{\lstinline!bcast\_if\_1!} and is available for explicit
use). When not specified, user provided dimensions get basis
\passthrough{\lstinline!default!}. This improves shape inference as it
ensures data vectors of dimension 1 don't become accidentally
broadcastable.

We track the distinction between parameter tensors and other leaf
tensors. When an axis of a parameter tensor has nothing constraining it,
inference fails with
\passthrough{\lstinline!"you forgot to specify a hidden dimension"!}.
When it is of a non-paramenter tensor, the unconstrained axis closes
upward to \(1_\emptyset\).

\section{Contexts and explicit
compilation}\label{contexts-and-explicit-compilation}

OCANNL has an immutable context \passthrough{\lstinline!Context.t!} that
unifies compilation, device, and execution contexts. Root contexts
determine the device. Here are the most important operations on a
context -- the bindings argument of \passthrough{\lstinline!compile!}
contains the externally bound indexing symbols used in the computation:

\begin{lstlisting}[language=Caml]
val compile : t -> Assignments.comp -> Indexing.unit_bindings -> t * routine
val run : t -> routine -> t
val copy : src:t -> dst:t -> Tnode.t -> unit
val to_host : t -> Tnode.t -> Ndarray.t
val from_host : t -> Tnode.t -> Ndarray.t -> t
\end{lstlisting}

Contexts track dependencies (e.g.~whether a tensor node is initialized
before use) both at routine's compile time and at runtime.

OCANNL has multi-stage shape inference. Stage 1 runs as tensor
expressions are constructed. Calling
\passthrough{\lstinline!Context.compile!} triggers stages 2 and later
for constraints that are in flight (global store) at that time.

\section{Related work}\label{related-work}

OCANNL's shape inference design arose from the author's appreciation for
the readability of einsum notation in JAX, and familiarity with
constraint-based type inference combining parametricity and subtyping.
In particular, variable elimination, such as the closure computation in
\emph{Pottier, ``A Framework for Type Inference with Subtyping'', ICFP
1998}, inspired OCANNL's shape inference solver. Other related works
discussed below are not influences; we only offer a selection as an
exhaustive treatment would require a survey paper.

Dependent types enable tracking tensor shapes in tensor types; they are
included in restricted forms in ergonomic programming languages with
type inference. The DML system \emph{Xi, Pfenning, ``Dependent Types in
Practical Programming'', 1999} paved the way to GADTs and refinement
types. The most prominent current example is Futhark \emph{Henriksen,
Elsman, Bailly, ``Shape-Constrained Array Programming with
Size-Dependent Types'', 2023}. It has capable type inference, but no
rank-polymorphism. This contrasts with Qube \emph{Trojahner, Grelck,
``Dependently typed array programs don't go wrong'', 2009}. Qube has
parametric rank-polymorphism but no type inference.

An interesting recent work is Star \emph{Bachurski, Mycroft, and
Orchard, ``Structuring Arrays with Algebraic Shapes'', 2025}. It
interprets algebraic data types as types of tensor indices. This is
reminiscent of OCANNL's connection between shapes and projections. Star
supports structural rank-polymoprhism via subtyping: as in OCANNL,
Star's shapes form a lattice (but as in type systems, the lattice is
distributive, while OCANNL's lattice is not distributive). Star does not
aim at tracking array (dimension) sizes. Star does not have shape
inference in the following sense: it does not synthesize new algebraic
data types.

APL and J introduced the rank-polymorphic array-programming tradition:
functions consume cells and are lifted over frames of extra axes, as
formalized in Remora \emph{Slepak, Shivers, Manolios, ``The Semantics of
Rank Polymorphism'', 2019}. This is structural rank-polymorphism. A
similarity with OCANNL is that the computation semantics are derived at
function application site. In OCANNL, freshened constraints are
introduced when a function computes tensor values (at OCaml runtime),
these constraints are used both for shapes and for indexing (computation
semantics). Refined Remora \emph{Matviichuk, Shivers, ``Refined Remora:
Constraining Array Shapes'', 2026} adds SMT-checked refinements over
dimensions and shape sequences, bringing expressivity from below to
beyond what's available in OCANNL. However, refinements are not
inferred.

\emph{Vasilache, Zinenko, Theodoridis et. al.~``Tensor Comprehensions:
Framework-Agnostic High-Performance Machine Learning Abstractions'',
2018} is the closest surface precedent: index variables induce ranges,
right-hand-side-only indices induce reductions, affine subscripts
express convolution, and underconstrained ranges require explicit
annotations. Production systems propagate symbolic or partial shapes
through graphs or IRs. There is no rank-polymorphism. The notation is
not as concise as in OCANNL: numeric precisions need to always be
specified (in OCANNL they are inferred bidirectionally), and tensor
fixed rank templates need to be declared.

\section{Conclusions and future work}\label{conclusions-and-future-work}

OCANNL uniquely combines parametric rank-polymorphism and structural
rank-polymorphism, and offers powerful shape inference. It aims to cover
most practical shape inference challenges, barring inherent
incompleteness that would require backtracking.

The Appendix presents definitions and theorems showing termination and
soundness, and illustrating the degree of (in)completeness. A full
treatment with proofs is available at:
\href{https://ahrefs.github.io/ocannl/docs/pdfs/ocannl-formal-core-technical-report.pdf}{ahrefs.github.io/ocannl/docs/pdfs/ocannl-formal-core-technical-report.pdf}.

One problem to explore in OCANNL's shape inference is the current
semantics of dimension basis comparisons:
\(1_\emptyset \neq 1_\texttt{default}\). For example, it makes scalars
incompatible with dimension-1 data vectors as same-size arguments to
einsum-based operations. This problem has not yet surfaced in practice:
``it's a feature, not a bug''. Once practical examples suffer from this
incompatibility, they might motivate a more sophisticated approach
(e.g.~basis polymorphism).

OCANNL also intends to provide to the OCaml ecosystem a compilation path
for tensor computations across various GPU backends: Nvidia (CUDA),
Apple Silicon (Metal), AMD (HIP). The compiler performs inlining and
Common Subexpression Elimination on the lowered IR (a loop nest
language). Tiling for threadblocks and tensor cores, and further
optimizations, are ongoing/future work.

\textbf{Authorship:} all the content above was written entirely by the
human author, without any AI/LLM feedback. The appendix below and the
accompanying
\href{https://ahrefs.github.io/ocannl/docs/pdfs/ocannl-formal-core-technical-report.pdf}{technical
report} were created by Claude Fable 5 (interactively, Appendix trimmed
down for brevity) and GPT 5.5 (final compilation and proof gaps).

\section{Appendix: Shape and Projections inference: semantics and
correctness}\label{appendix-shape-and-projections-inference-semantics-and-correctness}

TR refers to
\href{https://ahrefs.github.io/ocannl/docs/pdfs/ocannl-formal-core-technical-report.pdf}{ahrefs.github.io/ocannl/docs/pdfs/ocannl-formal-core-technical-report.pdf}.

\subsection{1. Dimensions}\label{dimensions}

\textbf{Definition {[}TR Definition 1.1{]} (Dimensions).} Fix a set
\(\mathcal{B}\) of basis tags with a distinguished tag
\(\texttt{default} \in \mathcal{B}\). The set of dimensions is
\[D \;=\; \{1_\emptyset\} \;\cup\; \{\, n_b \mid n \ge 1,\; b \in \mathcal{B} \,\}.\]
Elements \(n_b\) are \emph{atoms}; \(1_\emptyset\) is the
\emph{claim-free unit}. Adjoin a fresh element \(\bot\) and set
\(D_\bot = D \cup \{\bot\}\).

\textbf{Definition {[}TR Definition 1.2{]} (Broadcast order).} On \(D\):
\(d_1 \sqsubseteq d_2\) iff \(d_2 = 1_\emptyset\) or \(d_1 = d_2\).
Extend to \(D_\bot\) by \(\bot \sqsubseteq d\) for all \(d\).

\textbf{Proposition {[}TR Proposition 1.3{]} (Lattice structure).}
\((D_\bot, \sqsubseteq)\) is a bounded lattice of height \(2\): top
\(1_\emptyset\), bottom \(\bot\), and the atoms form an antichain, each
covering \(\bot\) and covered by \(1_\emptyset\). Meets and joins:
\[d \wedge d = d,\quad d \wedge 1_\emptyset = d,\quad d_1 \wedge d_2 = \bot \ (d_1 \ne d_2 \text{ atoms});\qquad d \vee d = d,\quad d \vee \bot = d,\quad d_1 \vee d_2 = 1_\emptyset \ (d_1 \ne d_2 \text{ atoms}).\]

\textbf{Proposition {[}TR Proposition 1.4{]} (Non-distributivity).}
\(D\) contains at least three atoms,
e.g.~\(\{1_{\texttt{default}}, 2_{\texttt{default}}, 3_{\texttt{default}}\} \subseteq D\),
so \(D_\bot\) is not distributive, and not modular.

\subsection{2. Rows}\label{rows}

\textbf{Definition {[}TR Definition 2.1{]} (Rows).} A \emph{row} is a
triple \(l \cdot \diamond \cdot r\) with \(l, r \in D^{*}\) (the
\emph{leading} and \emph{trailing flanks}).
\(\mathrm{rank}(l \cdot \diamond \cdot r) = |l| + |r|\). Equality of
rows is componentwise equality of \((l, r)\) --- in particular
\textbf{marker-sensitive}:
\([3]\cdot\diamond\cdot[4] \ne [3,4]\cdot\diamond\cdot[\,]\). (This is
identity of the algebra's elements; it is \emph{not} the satisfaction
relation of equality \emph{constraints}, which TR Definition 3.4 judges
by the flat equivalence \(\approx\) of TR Definition 2.8.)

\textbf{Definition {[}TR Definition 2.2{]} (Expansion).} For a row
\(R = l \cdot \diamond \cdot r\) and \(n \ge \mathrm{rank}(R)\), the
expansion \(R{\uparrow}n \in D^{n}\) is the flat sequence
\(l \cdot (1_\emptyset)^{\,n - \mathrm{rank}(R)} \cdot r\).

\textbf{Definition {[}TR Definition 2.3{]} (Row order).}
\(l_2 \cdot \diamond \cdot r_2 \;\sqsubseteq\; l_1 \cdot \diamond \cdot r_1\)
iff (i) \(|l_1| \le |l_2|\) and \(|r_1| \le |r_2|\) (\emph{flanks fit};
note this implies \(\mathrm{rank}_2 \ge \mathrm{rank}_1\)), and (ii)
\((l_2 \cdot r_2)[i] \sqsubseteq (l_1\cdot\diamond\cdot r_1){\uparrow}\mathrm{rank}_2\,[i]\)
for all \(i\) (\emph{pointwise refinement of the expansion}).

\textbf{Lemma {[}TR Lemma 2.4{]} (Expansion monotonicity).} If
\(R_2 \sqsubseteq R_1\) then for every \(n \ge \mathrm{rank}(R_2)\):
\(R_2{\uparrow}n\,[i] \sqsubseteq R_1{\uparrow}n\,[i]\) for all \(i\).

\textbf{Proposition {[}TR Proposition 2.5{]} (Partial order).}
\(\sqsubseteq\) is a partial order on rows.

\textbf{Proposition {[}TR Proposition 2.6{]} (The empty row is the
greatest element).} For every row \(R\):
\(R \sqsubseteq [\,]\cdot\diamond\cdot[\,]\), and nothing else is above
all rows.

\textbf{Proposition {[}TR Proposition 2.7{]} (Rows form a
join-semilattice with top).} Any two rows \(R_1, R_2\) have a least
upper bound:
\[R_1 \vee R_2 \;=\; l_\vee \cdot \diamond \cdot r_\vee, \qquad |l_\vee| = \min(|l_1|,|l_2|),\ \ l_\vee[i] = l_1[i] \vee l_2[i],\]
and symmetrically for \(r_\vee\) aligned from the back (joins taken in
\(D\); note two distinct atoms join to \(1_\emptyset\), so
\(l_\vee \in D^{*}\)).

\textbf{Remark {[}TR §2{]} (Rank-polymorphism lives in the order).} TR
Definition 2.3 relates rows of different ranks directly; no quantifier
appears anywhere in this section.

\textbf{Definition {[}TR Definition 2.8{]} (Flat equivalence).}
\(R \approx R'\) iff \(\mathrm{flat}(R) = \mathrm{flat}(R')\) as
sequences, where \(\mathrm{flat}(l\cdot\diamond\cdot r) = l \cdot r\)
--- equivalently, \(\mathrm{rank}(R) = \mathrm{rank}(R')\) and
\(R{\uparrow}\mathrm{rank}(R) = R'{\uparrow}\mathrm{rank}(R')\), since a
row's expansion at its own rank is its flattening (TR Definition 2.2
with no padding). \(\approx\) is the marker-erasing equivalence; it is
what the implementation's closed-row ``equality'' computes (§10).

\textbf{Proposition {[}TR Proposition 2.9{]} (\(\approx\) versus the
order).} (i) The equivalence induced by \(\sqsubseteq\) (mutual
inequality) is the \textbf{identity} --- this is just antisymmetry (TR
Proposition 2.5). So \(\approx\) is not derivable from the order. (ii)
Stronger: any two \emph{distinct} \(\approx\)-related rows are
\(\sqsubseteq\)-\textbf{incomparable}. (iii) \(\approx\) is not a
\(\sqsubseteq\)-congruence; indeed no nontrivial equivalence is
compatible with the order. (iv) \(\approx\) is the kernel of expansion
at own rank; everything else the order consults --- flank-fit, and
expansions at strictly higher ranks (broadcasting) --- is exactly where
\(\approx\)-related rows diverge.

The implementation does not use the observational equivalence of
rows-as-operands: its equality sites use \(\approx\) (observational
equivalence is finer), its order sites the marked order (observational
equivalence is looser because it absorbs explicit inner-edge
\(1_\emptyset\) entries into the middle).

\textbf{Remark {[}TR after Proposition 2.9; §6.1{]}
(Equality-as-\(\approx\) is underdetermined; the marker placement is a
policy choice).} Reading the ground meaning of closed-row equality as
\(\approx\) (forced by practice: an einsum spec row must equate with a
literal shape whose marker sits at the front) makes row-equality
constraints underdetermined: in
\(l_1\cdot\langle\rho\rangle\cdot r_1 \approx C\), the flat content of
\(\gamma\rho\) is forced (the middle of \(\mathrm{flat}(C)\)) but its
marker placement is free, and by TR Proposition 2.9(ii) the candidates
are pairwise \(\sqsubseteq\)-incomparable --- a later \emph{inequality}
on \(\rho\) can distinguish them. Witness:
\(\Phi = \{\,\langle\rho\rangle \approx [\,]\cdot\diamond\cdot[3,5],\ \ [3]\cdot\diamond\cdot[9,5] \sqsubseteq \langle\rho\rangle\,\}\).
Under \(\approx\), \(\gamma\rho = [3]\cdot\diamond\cdot[5]\) satisfies
both (\([3,9,5]\) vs \([3,1_\emptyset,5]\)). The implementation is
\emph{\(\approx\)-checking but marked-committing}: it accepts the flat
alignment, then commits the closed side's structural split as the
value's marker placement --- on this \(\Phi\) it fails. TR Definition
3.4 adopts the \(\approx\) reading, so that failure is officially a
\emph{policy rejection} (TR Lemma 5.1's policy steps), not semantic
unsatisfiability: closed-row equality joins closing (the
non-principality remark below) as a deliberately non-principal site. A
potential improvement in coverage would be to delay row equivalence
constraints till a later stage of solving, but that would have its own
trade-offs.

\textbf{Definition {[}TR Definition 2.10{]} (Two-sorted ground rows).}
Adjoin to the marked rows of TR Definition 2.1 a second sort of
\emph{rigid} rows: \(F^\bullet\) with \(F \in D^n\), a flat sequence
with \textbf{no marker}; \(\mathrm{rank}(F^\bullet) = n\), and the
expansion \(F^\bullet{\uparrow}m\) is defined only at \(m = n\), where
it is \(F\). Extend \(\sqsubseteq\): (a) marked--marked: TR Definition
2.3 unchanged; (b) \(F_2^\bullet \sqsubseteq F_1^\bullet\) iff equal
rank and pointwise refinement; (c) \(F^\bullet \sqsubseteq R\) (\(R\)
marked) iff \(\mathrm{rank}(F) \ge \mathrm{rank}(R)\) and
\(F[i] \sqsubseteq R{\uparrow}\mathrm{rank}(F)\,[i]\) pointwise --- a
rigid result may refine a broadcastable operand; (d)
\(R \sqsubseteq F^\bullet\): \textbf{never} --- nothing broadcastable
sits below a rigid row.

\textbf{Proposition {[}TR Proposition 2.11{]} (The two-sorted order).}
The extended relation is a partial order; the empty marked row remains
the unique top; and on the rigid sort, flat equality is trivially a
congruence, since no context consults a marker that does not exist.
Admitting even rank-equal \(R \sqsubseteq F^\bullet\) would break
transitivity:
\([5]\cdot\diamond\cdot[3] \sqsubseteq [\,]\cdot\diamond\cdot[3] \sqsubseteq [3]^\bullet\)
would demand \([5]\cdot\diamond\cdot[3] \sqsubseteq [3]^\bullet\), a
rank mismatch against a rigid row. The desired relation exists in the
system as the derived \emph{one-shot} relation,
``\(R \sqsubseteq F^\bullet\) at equal rank, pointwise'', see
\(R \sqsubseteq^{1} F^\bullet\) of TR Definition 2.12.

\subsection{3. Terms, substitutions, and constraint
derivation}\label{terms-substitutions-and-constraint-derivation}

\textbf{Definition {[}TR Definition 3.1{]} (Terms).} Dimension terms
\(t ::= \alpha \mid n_b \mid 1_\emptyset\). Row terms
\(R ::= l \cdot \diamond \cdot r \mid l \cdot \langle\rho\rangle \cdot r\)
with \(l, r\) sequences of dimension terms; we identify a bare row
variable \(\rho\) with \([\,]\cdot\langle\rho\rangle\cdot[\,]\). A term
is \emph{ground} if variable-free. Each row term contains \textbf{at
most one} row variable, at the marker --- an invariant preserved by
everything below.

\textbf{Definition {[}TR Definition 3.2{]} (Substitution).} A
substitution \(\sigma\) is a finite, sort-respecting map from variables
to terms with
\(\mathrm{dom}(\sigma) \cap \mathrm{vars}(\mathrm{ran}(\sigma)) = \emptyset\)
(idempotency). Application is structural except at the marker: if
\(\sigma(\rho) = l' \cdot \langle\rho'\rangle \cdot r'\) (or closed,
\(l'\cdot\diamond\cdot r'\)), then
\[\sigma(l \cdot \langle\rho\rangle \cdot r) \;=\; (\sigma l \cdot l') \cdot \langle\rho'\rangle \cdot (r' \cdot \sigma r) \quad\big(\text{resp. } (\sigma l \cdot l') \cdot \diamond \cdot (r' \cdot \sigma r)\big).\]

\textbf{Lemma {[}TR Lemma 3.3{]} (Composition).}
\((\sigma \circ \tau)(X) = \sigma(\tau(X))\) for all terms \(X\), where
\(\sigma \circ \tau\) is the usual composite substitution.

\textbf{Definition {[}TR Definitions 3.4-3.5{]} (Semantics).} A
\emph{ground substitution} \(\gamma\) is total on a fixed countable
variable universe and maps into ground dimensions/rows. For a constraint
set \(\Phi\) over atomic constraints
\(t_1 = t_2 \mid t_1 \sqsubseteq t_2 \mid R_1 \approx R_2 \mid R_1 \sqsubseteq R_2\):
\[\mathrm{Sol}(\Phi) = \{\gamma \text{ ground} \mid \gamma(\phi) \text{ holds for every } \phi \in \Phi\},\]
where ground \emph{dimension} equalities are identity, ground \emph{row}
equalities are judged by the flat equivalence \(\approx\) of TR
Definition 2.8, equivalently, identity of the rigidifications
\(\mathrm{flat}(\cdot)^\bullet\) in the two-sorted algebra of TR
Definition 2.10, equivalently mutual \(\sqsubseteq^{1}\); the three
formulations coincide. A substitution \(\sigma\) \emph{models} \(\Phi\),
\(\sigma \models \Phi\), iff
\(\gamma \circ \sigma \in \mathrm{Sol}(\Phi)\) for every ground
\(\gamma\), equivalently, \(\sigma(\Phi)\) is valid under universal
quantification of its free variables. The \emph{substitution order}:
\(\sigma_1 \le \sigma_2\) iff
\(\exists u.\ u \circ \sigma_1 = \sigma_2\).

\textbf{Example {[}TR Example 3.6{]} (No principal model exists in
general).} Let \(\Phi = \{3_b \sqsubseteq \alpha\}\). The identity does
not model \(\Phi\) (ground \(\alpha \mapsto 5_b\) falsifies it). Any
model must map \(\alpha\) to a term all of whose groundings lie above
\(3_b\), i.e.~to \(3_b\) or \(1_\emptyset\) (a variable target fails as
the identity did). The two models \(\sigma_1 = [\alpha \mapsto 3_b]\)
and \(\sigma_2 = [\alpha \mapsto 1_\emptyset]\) are
\(\le\)-incomparable: a witness \(u\) for either direction would have to
send a ground dimension to a different ground dimension, which
substitutions cannot do. Hence \(\Phi\) has models but \textbf{no
\(\le\)-least model}.

\emph{Consequence.} The solver's answer is a pair --- a substitution
\emph{and a residual bound store} --- and the correct theorem is about
solution-set representation and most-generality (TR Theorem 5.9 below).

\textbf{Compact rule table (constraint derivation).} Let \(C\) be the
result shape and \(A,B,D\) operand shapes, with \(S_k\) denoting the row
of kind \(k \in \{B,I,O\}\) (batch, input, output). For spec-based
rules, let \(T^{\mathrm{lhs}}_k\) and \(T^{\mathrm{rhs}}_{i,k}\) be the
rows obtained by elaborating the result and operand slots of the
specification. The generated core constraints are:

\begin{longtable}[]{@{}
  >{\raggedright\arraybackslash}p{(\linewidth - 2\tabcolsep) * \real{0.5000}}
  >{\raggedright\arraybackslash}p{(\linewidth - 2\tabcolsep) * \real{0.5000}}@{}}
\toprule\noalign{}
\begin{minipage}[b]{\linewidth}\raggedright
Operation logic
\end{minipage} & \begin{minipage}[b]{\linewidth}\raggedright
Generated constraints
\end{minipage} \\
\midrule\noalign{}
\endhead
\bottomrule\noalign{}
\endlastfoot
transpose &
\(C_B \sqsubseteq A_B,\quad C_I \sqsubseteq A_O,\quad C_O \sqsubseteq A_I\) \\
pointwise unary & \(C_k \sqsubseteq A_k\) for every \(k\) \\
pointwise binary & \(C_k \sqsubseteq A_k,\quad C_k \sqsubseteq B_k\) for
every \(k\) \\
pointwise ternary &
\(C_k \sqsubseteq A_k,\quad C_k \sqsubseteq B_k,\quad C_k \sqsubseteq D_k\)
for every \(k\) \\
compose &
\(A_I \sqsubseteq B_O,\quad C_B \sqsubseteq A_B,\quad C_B \sqsubseteq B_B,\quad C_I \sqsubseteq B_I,\quad C_O \sqsubseteq A_O\) \\
compose accumulate & the compose constraints, plus
\(C_k \sqsubseteq D_k\) for every \(k\) \\
batch slice &
\((s \cdot C_B) \approx A_B,\quad C_I \approx A_I,\quad C_O \approx A_O\),
where \(s\) is the sliced outer batch axis \\
permute &
\(C_k \approx T^{\mathrm{lhs}}_k,\quad T^{\mathrm{rhs}}_{1,k} \approx A_k\)
for every \(k\) \\
einsum &
\(C_k \approx T^{\mathrm{lhs}}_k,\quad T^{\mathrm{rhs}}_{i,k} \approx A_{i,k}\)
for every operand \(i\) and kind \(k\) \\
\end{longtable}

Here \(\sqsubseteq\) is the broadcast/refinement order, while
\(\approx\) is flat row equivalence. Thus pointwise, transpose, and
compose operations generate the broadcast-checking constraints of the
core; permute and einsum deliberately generate equalities, the stronger
inference policy used to recover labels and row variables
bidirectionally.

\subsection{4. The solver: configurations and
rules}\label{the-solver-configurations-and-rules}

\textbf{Definition {[}TR Definition 4.1{]} (Configuration).} A
configuration is \(\langle \Phi;\ \sigma;\ B \rangle\): unsolved
constraints, accumulated (idempotent) bindings, and a \emph{bound
store}. \(B\) assigns to each unsolved variable: for a dimension
variable \(\alpha\), at most one atom lower bound
\(\mathrm{lb}(\alpha) \in \{n_b\}\) plus a set of deferred
variable--variable inequalities (``adjacencies''); for a row variable
\(\rho\), a set of row-term lower bounds (``caps'', from inequality
residues) plus adjacencies. There is a distinguished failure
configuration \(\mathsf{fail}\).

Binding a variable (\(\sigma := [x \mapsto T] \circ \sigma\))
\textbf{re-emits} into \(\Phi\) every bound and adjacency stored for
\(x\), instantiated at \(T\), and substitutes \(T\) for \(x\) throughout
\(\Phi\) and \(B\). This single discipline replaces ad-hoc propagation.

At dimension sort, the stored \(\mathrm{glb}(\alpha)\) is understood as
the join of all ground lower bounds known for \(\alpha\), including
those entailed through directed adjacencies. Operationally, when
\(d \sqsubseteq \alpha\) is recorded, the solver enqueues
\(d \sqsubseteq \beta\) for every stored edge
\(\alpha \sqsubseteq \beta\); when a new edge
\(\alpha \sqsubseteq \beta\) is recorded, it enqueues every stored cap
\(d \sqsubseteq \alpha\) as \(d \sqsubseteq \beta\). Fair processing
iterates these emissions to fixpoint. This is not a closing policy: it
is the transitivity step
\(d \sqsubseteq \alpha \sqsubseteq \beta \Rightarrow d \sqsubseteq \beta\),
and DI-cap handles any collapse produced downstream.

\textbf{Definition {[}TR §4.1; TR Lemma 4.2{]} (Dimension rules).}

\begin{longtable}[]{@{}
  >{\raggedright\arraybackslash}p{(\linewidth - 4\tabcolsep) * \real{0.3333}}
  >{\raggedright\arraybackslash}p{(\linewidth - 4\tabcolsep) * \real{0.3333}}
  >{\raggedright\arraybackslash}p{(\linewidth - 4\tabcolsep) * \real{0.3333}}@{}}
\toprule\noalign{}
\begin{minipage}[b]{\linewidth}\raggedright
\end{minipage} & \begin{minipage}[b]{\linewidth}\raggedright
Constraint
\end{minipage} & \begin{minipage}[b]{\linewidth}\raggedright
Action
\end{minipage} \\
\midrule\noalign{}
\endhead
\bottomrule\noalign{}
\endlastfoot
DE-refl & \(t = t\) & discard \\
DE-clash & \(n_b = m_c\), \((n,b)\ne(m,c)\) (incl.~atom vs
\(1_\emptyset\)) & \(\mathsf{fail}\) \\
DE-bind & \(\alpha = t\), \(t \ne \alpha\) & bind
\(\alpha \mapsto t\) \\
DI-top & \(t \sqsubseteq 1_\emptyset\) & discard \\
DI-refl & \(t \sqsubseteq t\) & discard \\
DI-ground & \(d \sqsubseteq d'\) ground & check TR Definition 1.2; else
\(\mathsf{fail}\) \\
DI-pin & \(\alpha \sqsubseteq d\), \(d\) an atom & replace by
\(\alpha = d\) \\
DI-pin-top & \(1_\emptyset \sqsubseteq \alpha\) & replace by
\(\alpha = 1_\emptyset\) \\
DI-cap & \(d \sqsubseteq \alpha\), \(d\) an atom & if
\(\mathrm{lb}(\alpha)\) unset: record and forward along stored outgoing
adjacencies; if \(= d\): discard; if \(= d' \ne d\): replace by
\(\alpha = 1_\emptyset\) \\
DI-adj & \(\alpha \sqsubseteq \beta\), \(\alpha \ne \beta\) & defer
(record adjacency on both) and forward \(\alpha\)'s stored cap, if any,
to \(\beta\) \\
\end{longtable}

\emph{Justifications (each rule preserves \(\mathrm{Sol}\)):} DI-pin: in
\(D\), the only element \(\sqsubseteq\) an atom \(d\) is \(d\) itself.
DI-pin-top: the only element above the top is the top. DI-cap collapse:
any \(\gamma(\alpha)\) above two distinct atoms must be \(1_\emptyset\)
(TR Proposition 1.3). DI-cap/DI-adj forwarding emits only transitive
consequences of already-stored constraints. DI-pin-top is the case the
blog's narrative omits.

\textbf{Definition {[}TR Definitions 4.3-4.4; TR Lemma 4.5{]} (Row
rules).} Write both sides aligned at the \textbf{outer edges}: leading
flanks from the front, trailing flanks from the back.

For equality, first emit dimension equalities on the overlapping aligned
flank positions. Let \(m_l, m_r\) be the surpluses: the unmatched inner
segments of the longer leading and trailing flanks. ``Defer into the
closing policy'' means postpone a constraint till a later stage of the
solving process.

\emph{Equality \(R_1 \approx R_2\).}

\begin{longtable}[]{@{}
  >{\raggedright\arraybackslash}p{(\linewidth - 4\tabcolsep) * \real{0.3333}}
  >{\raggedright\arraybackslash}p{(\linewidth - 4\tabcolsep) * \real{0.3333}}
  >{\raggedright\arraybackslash}p{(\linewidth - 4\tabcolsep) * \real{0.3333}}@{}}
\toprule\noalign{}
\begin{minipage}[b]{\linewidth}\raggedright
\end{minipage} & \begin{minipage}[b]{\linewidth}\raggedright
Situation
\end{minipage} & \begin{minipage}[b]{\linewidth}\raggedright
Action
\end{minipage} \\
\midrule\noalign{}
\endhead
\bottomrule\noalign{}
\endlastfoot
RE-closed & both sides closed & \(\mathsf{fail}\) iff the ranks differ;
otherwise zip the flattened rows, emitting
\(\mathrm{flat}(R_1)[i] = \mathrm{flat}(R_2)[i]\) \\
RE-open-closed & one side open with \(\rho\), the other closed &
\(\mathsf{fail}\) iff the open side's total explicit flank length
exceeds the closed side's rank; otherwise bind \(\rho\) to the uncovered
middle of \(\mathrm{flat}(R_{\mathrm{closed}})\), placed by the
inherit-the-split policy \\
RE-nested & both sides open, \(\rho_1 \ne \rho_2\), and both surpluses
are on \(\rho_2\)'s side & bind
\(\rho_1 \mapsto m_l \cdot \langle\rho_2\rangle \cdot m_r\) \\
RE-cross & both sides open, \(\rho_1 \ne \rho_2\), with split surpluses
& defer the residual word equation \(s \cdot x_1 = x_2 \cdot t\) into
the closing policy \\
RE-same-empty & both sides open with the same \(\rho\) and no surplus &
discard \\
RE-same-occurs & both sides open with the same \(\rho\) and unequal
known flank totals & \(\mathsf{fail}\) \\
RE-same-rot & both sides open with the same \(\rho\) and equal totals
but shifted splits & defer the rotational equation
\(x \cdot t = s \cdot x\) into the closing policy \\
\end{longtable}

\emph{Justifications and policy choices.} RE-closed is exactly TR
Definition 3.4's \(\approx\)-semantics: markers are not consulted. In
RE-open-closed, substitution extends the open side's flanks inward only,
so no grounding can shorten them. The binding
\(\rho \mapsto m_l \cdot \diamond \cdot m_r\) with \(m_l \cdot m_r = m\)
additionally chooses a placement that \(\approx\) leaves free (the
marker-placement remark above): it commits to the closed side's declared
split where that split falls inside \(m\), and to the left edge
otherwise. This is a policy step in the family of TR Definition 6.1.
\emph{(The implementation additionally rejects when the open side's
trailing flank overflows the closed side's; conservative under
\(\approx\).)}

RE-nested is entailed by cancelling common outer flanks of the flat
words; only placement is policy, as above. RE-cross is the two-variable
word equation case: its solutions are the principal family
\(x_2 = s\cdot w\), \(x_1 = w\cdot t\) \textbf{plus} at most \(|s|\)
\emph{sporadic} cross-overlap solutions (\(x_2\) a proper prefix of
\(s\), the rest forced). Consider an alternative: fresh-variable binding
\(\rho_2 \mapsto s\cdot\langle\rho'\rangle\),
\(\rho_1 \mapsto \langle\rho'\rangle\cdot t\). It captures exactly the
principal family: exact under the former marked semantics, but under
\(\approx\) it forecloses the sporadic solutions, a \emph{rank}
commitment stronger than a placement choice (and it breaks rank
conservativity of TR Lemma 5.3). Deferring keeps the equation in flight;
substituting either variable makes the check exact, and the closing
stage's upward close selects the corresponding extremal solution (for
the variable carrying \(x_2\), the outermost \emph{sporadic} one), then
the re-emitted equation validates it.

RE-same-occurs fails because no finite row satisfies
\(\rho \approx m_l \cdot \langle\rho\rangle \cdot m_r\) with
\(m_l \cdot m_r\) nonempty. RE-same-rot is different. Its residue
\(x \cdot t = s \cdot x\) is satisfiable under TR Definition 3.4's
\(\approx\)-semantics exactly for conjugate \(s, t\) with cyclically
periodic \(x\) (Lyndon--Schützenberger), a solution family outside the
solver's constraint language. Deferring lets other constraints solve
\(\rho\) and make the substituted closed-closed flat check exact.
Otherwise, guessing closes \(\rho\) upward --- the least-material
disjunct, \(x = [\,]\) --- after which the equation requires \(s = t\).
This accepts exactly the \(\delta = 0\) conjugacy class: a policy
commitment in the family of TR Definition 6.1 and since \(\approx\) is
the adopted semantics, a genuine incompleteness on nontrivial conjugate
instances (periodic \(x \ne [\,]\)), accepted deliberately.

For inequality, first emit dimension inequalities
\(t_{\mathrm{res}} \sqsubseteq t_{\mathrm{op}}\) on aligned flank
overlaps.

\emph{Inequality \(R_{\mathrm{res}} \sqsubseteq R_{\mathrm{op}}\).}

\begin{longtable}[]{@{}
  >{\raggedright\arraybackslash}p{(\linewidth - 4\tabcolsep) * \real{0.3333}}
  >{\raggedright\arraybackslash}p{(\linewidth - 4\tabcolsep) * \real{0.3333}}
  >{\raggedright\arraybackslash}p{(\linewidth - 4\tabcolsep) * \real{0.3333}}@{}}
\toprule\noalign{}
\begin{minipage}[b]{\linewidth}\raggedright
\end{minipage} & \begin{minipage}[b]{\linewidth}\raggedright
Situation
\end{minipage} & \begin{minipage}[b]{\linewidth}\raggedright
Action
\end{minipage} \\
\midrule\noalign{}
\endhead
\bottomrule\noalign{}
\endlastfoot
RI-cap & operand open, and result supplies at least the operand's flanks
& record the result's interior residue as a \textbf{cap} on
\(\rho_{\mathrm{op}}\); the residue is a row term, open if the result is
open \\
RI-closed-op & operand closed, and result rank is at least the operand
flanks & discard the operand's interior expansion; it is
\(1_\emptyset\), and the corresponding result positions are
unconstrained \\
RI-short-closed & result closed and shorter than the operand's known
flanks & \(\mathsf{fail}\) \\
RI-deficit & result open and shorter than the operand's known flanks,
with deficit \(k > 0\) & bind
\(\rho_{\mathrm{res}} \mapsto \alpha_1 \cdots \alpha_{k_l} \cdot \langle\rho'\rangle \cdot \beta_1 \cdots \beta_{k_r}\)
with fresh dimension variables and a fresh row variable, sized to clear
the deficit; then reprocess \\
\end{longtable}

RI-short-closed is a genuine rank mismatch. RI-deficit is subject to the
rank-cycle check of TR Definition 4.6.

\textbf{Definition {[}TR §5{]} (Fairness).} A run processes constraints
until \(\Phi = \emptyset\) or \(\mathsf{fail}\); re-emitted constraints
return to \(\Phi\). Any fair strategy is allowed.

\textbf{Example {[}TR §4.3{]} (Divergence without a cycle check).} Let
\(\Phi_0 = \{\, \rho_1 \sqsubseteq [a]\cdot\langle\rho_2\rangle,\ \ \rho_2 \sqsubseteq [b]\cdot\langle\rho_1\rangle \,\}\)
(\(a, b\) ground atoms). Semantically: the first constraint forces
\(\mathrm{rank}(\gamma\rho_1) \ge 1 + \mathrm{rank}(\gamma\rho_2)\)
(flanks-fit in TR Definition 2.3), the second forces
\(\mathrm{rank}(\gamma\rho_2) \ge 1 + \mathrm{rank}(\gamma\rho_1)\);
hence \(\mathrm{Sol}(\Phi_0) = \emptyset\). Operationally, RI-deficit
fires on \(\rho_1\), growing it by one axis; substitution lengthens the
\emph{operand} side of the second constraint, where RI-deficit grows
\(\rho_2\); substitution lengthens the operand side of the first
constraint's residue, and so on forever. The dimension and row rules
alone do not terminate. The self-cyclic instance
\(\rho \sqsubseteq [a]\cdot\langle\rho\rangle\) (obtainable by an einsum
equality merging two row variables) diverges the same way. Developing
this formalization uncovered an implementation bug: due to other
mechanisms, the minimal \emph{live} divergence was the three-variable
cycle
\(\{\rho_1 \sqsubseteq \langle\rho_2\rangle{\cdot}[a],\ \rho_2 \sqsubseteq \langle\rho_3\rangle{\cdot}[b],\ \rho_3 \sqsubseteq \langle\rho_1\rangle{\cdot}[c]\}\),
which ran unboundedly until killed (now rejected by the TR Definition
4.6 check).

\textbf{Definition {[}TR Definition 4.6{]} (Rank-fact graph and cycle
check).} Maintain beside the configuration a directed graph \(G\) on row
variables --- \textbf{persistent}: edges are only ever added, never
retracted, in particular not when a variable is solved and substituted
away. An edge \(\rho \xrightarrow{\,k\,} \rho'\) with weight \(k \ge 0\)
records the entailed fact
\(\mathrm{rank}(\gamma\rho) \ge \mathrm{rank}(\gamma\rho') + k\) (for
every solution \(\gamma\) of the current configuration). Edges are
recorded at three points:

\begin{enumerate}
\def\labelenumi{(\roman{enumi})}
\tightlist
\item
  \emph{binding:} solving
  \(\rho \mapsto l \cdot \langle\rho'\rangle \cdot r\) records
  \(\rho \xrightarrow{\,|l|+|r|\,} \rho'\) (the entailed fact is an
  equality; one direction suffices for the check; this covers RE-nested
  bindings and open row bindings produced by RI-deficit);
\item
  \emph{equal flanks:} processing
  \(R_{\mathrm{res}} \sqsubseteq R_{\mathrm{op}}\) with both sides open
  and equal known flank lengths records
  \(\rho_{\mathrm{res}} \xrightarrow{\,0\,} \rho_{\mathrm{op}}\);
\item
  \emph{deficit:} RI-deficit with deficit \(k > 0\) and open operand
  records \(\rho_{\mathrm{res}} \xrightarrow{\,k\,} \rho_{\mathrm{op}}\)
  \emph{before} growing (a closed operand needs no edge: the rank bound
  is absolute and the growth is one-shot).
\end{enumerate}

Every insertion is \textbf{guarded}: if the new edge closes a directed
cycle of total weight \(> 0\), then \(\mathsf{fail}\).

\emph{Soundness of failing:} summing the facts around a positive cycle
gives \(\mathrm{rank}(\gamma\rho) \ge \mathrm{rank}(\gamma\rho) + w\)
with \(w > 0\), so \(\mathrm{Sol} = \emptyset\). \emph{Soundness of
persistence:} each edge is entailed by constraints present when it was
recorded, and the solver never retracts a constraint --- TR Lemma 5.1
only moves and re-expresses them --- so entailed facts stay entailed,
even when the variables involved have long been substituted away.
Zero-weight cycles are legal: they assert rank \emph{equality} along the
cycle (and could be collapsed to row equalities, generalizing the
one-step row antisymmetry the implementation already performs ---
cf.~§10).

\emph{Sign discipline.} In RI-cap --- operand open, result's known
flanks in \emph{surplus} by \(s > 0\) --- the constraint entails only
\(\mathrm{rank}(\gamma\rho_{\mathrm{res}}) \ge \mathrm{rank}(\gamma\rho_{\mathrm{op}}) - s\):
a \emph{nonpositive} lower bound. Recording it with weight \(0\) (or
\(s\)) is \textbf{unsound}: it manufactures cycles that no solution
violates. A first implementation attempt did exactly this and rejected
legitimate programs (the SGD self-reference idiom among them); the fixed
implementation records nothing there. Recording a genuinely negative
weight would be sound but is unnecessary for the check's purpose.

\textbf{Proposition {[}TR Proposition 4.7{]} (The check is exact for the
recorded facts).} At any point in a run, the recorded fact set
\(\{\mathrm{rank}(\rho) \ge \mathrm{rank}(\rho') + k\}\) is satisfiable
over \(\mathbb{N}\) iff \(G\) has no positive-weight cycle. Hence the
guard fails exactly when the facts recorded so far are jointly
unsatisfiable.

\emph{(Standard difference-constraints/Bellman--Ford duality --- worth
stating because it delimits what the check decides: unsatisfiability of
the \textbf{recorded} facts, a sound under-approximation of
rank-unsatisfiability of the configuration. Whether enough facts are
always recorded in time is exactly TR Lemma 5.5.)}

\subsection{5. Metatheory of solving}\label{metatheory-of-solving}

\textbf{Lemma {[}TR Lemma 5.1{]} (Solution preservation, with a policy
coordinate).} Write \(\widehat{C}\) for the constraint set
\(\Phi \cup \mathrm{eqns}(\sigma) \cup \mathrm{constr}(B)\) denoted by a
configuration (\(\mathrm{eqns}(\sigma) = \{x \approx \sigma(x)\}\);
\(\mathrm{constr}(B)\) the stored bounds/adjacencies as constraints),
and consider \(\mathrm{Sol}(\widehat{C})\) under TR Definition 3.4. The
rules above divide:

\emph{Semantic steps} --- all dimension rules, all row-inequality rules
(RI-cap, RI-closed-op, RI-short-closed, RI-deficit), and the checking
content of the equality rules (the flat zip of RE-closed, the outer-edge
dimension equalities, the rank-overflow failure of RE-open-closed, the
discard of RE-same-empty, the occurs failure of RE-same-occurs, and the
entailed flat content of RE-open-closed and RE-nested bindings) ---
preserve \(\mathrm{Sol}(\widehat{C})\) exactly, up to extension of
\(\gamma\) to freshly introduced variables (RI-deficit): every solution
of the old configuration extends to one of the new, and every solution
of the new restricts to one of the old. \(\mathsf{fail}\) steps of this
kind are taken only when \(\mathrm{Sol}(\widehat{C}) = \emptyset\).

\emph{Policy steps} --- the \textbf{placement} component of the
row-variable bindings in RE-open-closed and RE-nested, and the
\textbf{deferred-equation resolutions} at closing for RE-cross and
RE-same-rot (both word-equation residues resolved by the upward close)
--- \emph{refine} \(\mathrm{Sol}\): the bound variable's flat content is
entailed, but the committed split selects one of the pairwise
\(\sqsubseteq\)-incomparable representatives of TR Proposition 2.9(ii),
and substituting the committed value into stored \emph{marked
inequalities} may strengthen them. Every solution of the new
configuration restricts to a solution of the old; conversely a solution
of the old survives into the new after re-placing the bound variables'
markers, which is possible iff no marked inequality of \(\widehat{C}\)
discriminates the placement --- the witness in the marker-placement
remark above is the counterexample, and the order-dependence discussion
in TR §6.1 is its operational trace. A \(\mathsf{fail}\) arising from a
substituted placement, where another representative would have passed,
is therefore a \emph{policy rejection}, not semantic unsatisfiability.

\textbf{Definition {[}TR Definition 5.2{]} (Rank abstraction).} For a
configuration with constraint denotation \(\widehat{C}\), its \emph{rank
abstraction} \(\lfloor\widehat{C}\rfloor\) is the finite set of rank
facts read off one constraint at a time: a row constraint relating
\(l\cdot\langle\rho\rangle\cdot r\) to
\(l'\cdot\langle\rho'\rangle\cdot r'\) (as result/operand of
\(\sqsubseteq\), resp. as an equality) contributes
\(\mathrm{rank}(\rho) \ge \mathrm{rank}(\rho') + (|l'|+|r'|) - (|l|+|r|)\)
(both directions for an equality); open-vs-closed pairs contribute the
corresponding absolute bounds on \(\mathrm{rank}(\rho)\). A \emph{rank
model} is an assignment \(r : \mathrm{RowVars} \to \mathbb{N}\)
satisfying \(\lfloor\widehat{C}\rfloor\). Every
\(\gamma \in \mathrm{Sol}(\widehat{C})\) induces a rank model
\(\rho \mapsto \mathrm{rank}(\gamma\rho)\) (flanks-fit in TR Definition
2.3); the converse fails --- the abstraction forgets dimensions and
bases.

\textbf{Lemma {[}TR Lemma 5.3{]} (Rank conservativity).} Every
non-\(\mathsf{fail}\) rule maps a configuration with a rank model to one
with a rank model agreeing on the old variables (extended on fresh
ones). RI-deficit extends by
\(r(\rho') := r(\rho_{\mathrm{res}}) - k \ge 0\), justified by the
triggering constraint's own rank fact; everything else preserves the
model verbatim. \emph{(A.2's remark: the previously drafted RE-cross
fresh-variable binding would break this lemma --- its
\(\mathrm{rank}(\rho_2) \ge |s|\) commitment is not implied by the
consumed fact --- an independent reason for RE-cross's correction to
deferral.)} \(\square\)

\textbf{Theorem {[}TR Theorem 5.4; TR Lemma 5.5{]} (Termination).} Let
\(\Phi_0\) be finite.

\textbf{(a)} If \(\Phi_0\) has a rank model --- in particular whenever
\(\mathrm{Sol}(\Phi_0) \ne \emptyset\) --- then every fair run
terminates, with or without the cycle check; moreover the check never
fires on such inputs (its recorded facts are entailed, hence satisfied
by the rank model, hence positive-cycle-free by TR Proposition 4.7), so
guarding is semantically free.

\textbf{(b)} If \(\Phi_0\) has no rank model, then
\(\mathrm{Sol}(\Phi_0) = \emptyset\) and no run reaches solved form (TR
Proposition 5.7 plus TR Lemma 5.1 would otherwise produce a solution,
hence a rank model); every run either fails or diverges, and the cycle
check is what must convert divergence into failure. A small gap remains
in our proof that it always does. The proof is built around:

\textbf{Detection Lemma {[}TR Lemma 5.5{]}.} Every infinite run inserts,
at some finite stage, an edge closing a positive-weight cycle in \(G\).
(Hence, with the guard of TR Definition 4.6, infinite runs do not exist,
and rank-unsatisfiable inputs fail in finite time.)

\textbf{Definition {[}TR Definition 5.6{]} (Solved form).} A final
configuration \(\langle \emptyset; \sigma_\star; B_\star \rangle\) is in
\emph{solved form}: \(\sigma_\star\) idempotent; every variable in
\(B_\star\) unsolved; each dimension variable carries at most one atom
lower bound; all stored row caps and adjacencies are between unsolved
variables and contain no pending pins (else a rule would fire).

\textbf{Proposition {[}TR Proposition 5.7{]} (The residual store is
always satisfiable; \(\gamma_\uparrow\) is its greatest solution).
{[}proved{]}} Define \(\gamma_\uparrow\) on the unsolved variables:
\(\gamma_\uparrow(\alpha) = 1_\emptyset\),
\(\gamma_\uparrow(\rho) = [\,]\cdot\diamond\cdot[\,]\), extended through
\(\sigma_\star\) on solved variables
(\(\gamma_\uparrow(x) = \gamma_\uparrow(\sigma_\star(x))\)). Then (i)
\(\gamma_\uparrow \in \mathrm{Sol}(\mathrm{constr}(B_\star) \cup \mathrm{eqns}(\sigma_\star))\);
(ii) \(\gamma_\uparrow\) is the pointwise-greatest element of that
solution set: for every solution \(\gamma\) and every variable \(x\),
\(\gamma(x) \sqsubseteq \gamma_\uparrow(x)\).

\emph{(Scope notes. First: solved form here is taken with the
RE-same-rot deferrals discharged --- an in-flight deferred equation
\(x \cdot t = s \cdot x\) is satisfied by \(\gamma_\uparrow\) iff
\(s = t\), so a surviving \(s \ne t\) deferral postpones (i) to the
closing step that resolves or fails it; the success direction of TR
Corollary 5.8 is unaffected, since an end-to-end successful run has
discharged its deferrals. Second: (ii)'s greatestness is over the
solution set of the} final configuration, \emph{which encodes the run's
placement commitments; it does not by itself give greatestness over
\(\mathrm{Sol}(\Phi_0)\) --- that transfer is the content, and the
caveat, of TR Proposition 6.4.)}

\textbf{Corollary {[}TR Corollary 5.8{]} (Decision).} A fair run fails
iff the \emph{policy-strengthened} system is unsatisfiable ---
\(\Phi_0\) plus TR Lemma 5.1's policy commitments (placements from
RE-open-closed and RE-nested, and deferred-equation resolutions from
RE-cross and RE-same-rot). Success implies
\(\mathrm{Sol}(\Phi_0) \ne \emptyset\) outright (TR Proposition 5.7
composed back through the preservation chain --- the success direction
is unconditional). Failure implies \(\mathrm{Sol}(\Phi_0) = \emptyset\)
\emph{when the failing step is semantic}; a policy rejection can fail an
\(\approx\)-satisfiable input (the witness in the marker-placement
remark above). On inputs satisfiable \emph{and}
placement-undiscriminated, runs terminate by TR Theorem 5.4 and cannot
fail, so they decide positively. On unsatisfiable inputs, runs cannot
succeed; that they \emph{terminate} --- making the solver a decision
procedure rather than a semi-decision procedure --- is exactly TR Lemma
5.5.

\textbf{Theorem {[}TR Theorem 5.9{]} (Soundness, representation, most
generality).} Let a fair run on finite \(\Phi_0\) terminate in
\(\langle\emptyset; \sigma_\star; B_\star\rangle\). Then, with
\(V = \mathrm{vars}(\Phi_0)\):

\begin{enumerate}
\def\labelenumi{\arabic{enumi}.}
\tightlist
\item
  \emph{(Representation)} \(\gamma \in \mathrm{Sol}(\Phi_0)\) iff
  \(\gamma\) extends to a ground \(\hat\gamma\) (on the fresh variables)
  with \(\hat\gamma = \hat\gamma \circ \sigma_\star\) and
  \(\hat\gamma \in \mathrm{Sol}(\mathrm{constr}(B_\star))\).
\item
  \emph{(No guessing / entailment)} every binding \(x \mapsto T\) in
  \(\sigma_\star\) is conservative: every solution of \(\Phi_0\) extends
  to one satisfying \(x = T\).
\item
  \emph{(Most generality among models)} for every substitution
  \(\sigma \models \Phi_0\) there is \(u\) with
  \(u \circ \sigma_\star = \sigma\) on \(V\) --- indeed \(u = \sigma\)
  works: \(\sigma = \sigma \circ \sigma_\star\) on \(V\).
\end{enumerate}

\textbf{Caveat to the whole theorem (the \(\approx\)/policy weakening,
with TR Definition 3.4's adoption):} (1)'s ``iff'' and (2)'s ``every
solution extends'' hold exactly for runs of semantic steps; across TR
Lemma 5.1's \emph{policy steps} the forward direction holds after
re-placing the policy-committed markers, hence unconditionally only on
placement-undiscriminated inputs (the conjectured surface fragment ---
TR Corollary 5.8). (3) is up to \(\approx\) as the proof now states:
\(\sigma_\star\) is most general in content; its marker placements are
the policy's, not canonical. This is the theorem-level form of the
order-dependence caveat in TR §6.1.

\textbf{Caveat to (1)/(2):} the extension over fresh variables is
existential, exactly as in standard unification with fresh-variable
introduction; the answer is unique up to renaming of fresh variables.

\textbf{Corollary {[}TR Theorem 5.9(2){]} (Forced variables).}
\(\sigma_\star\) binds only variables whose values (relative to the
remaining free ones) are entailed by \(\Phi_0\) --- by TR Theorem
5.9(2). Conversely every variable left unsolved genuinely admits at
least two solutions or carries only its top default (\(\gamma_\uparrow\)
vs.~e.g.~committing a cap), by TR Proposition 5.7 and inspection of
solved form.

\textbf{Corollary {[}TR Corollary 5.10{]} (Principal model, recovered).}
If \(B_\star\) stores no atom caps and no row caps (adjacencies
allowed), then \(\sigma_\star\) (read as a model, free variables
universal) satisfies \(\sigma_\star \models \Phi_0\), and by TR Theorem
5.9(3) it is the \(\le\)-least model: the blog's ``principal model''
claim holds exactly in this case. TR Example 3.6 shows the restriction
is necessary.

\subsection{6. Closing}\label{closing}

\textbf{Definition {[}TR Definition 6.1{]} (Closing policy).} Variables
carry a provenance tag: \emph{leaf} (belonging to a terminal tensor's
shape; a subset are \emph{parameters}) or \emph{interior}. Closing runs
after solving, so dimension leaves close to the \(\mathrm{glb}\) defined
by the saturated store of TR Definition 4.1. Closing then proceeds
interleaved: 1. \emph{(Close leaves, downward.)} For each unsolved leaf
variable: bind a dimension variable to its stored atom cap if any, else
to \(1_\emptyset\) --- except a \textbf{parameter} with no cap at some
position: \(\mathsf{error}\) (``unspecified hidden dimension''). Bind a
row variable to the join (TR Proposition 2.7) of the \emph{ground parts}
of its stored caps, with \emph{holes} (positions where no cap is ground)
filled by \(1_\emptyset\), and no-further-axes beyond the caps' extent;
parameters error at holes. 2. \emph{(Re-solve.)} The bindings re-emit
the variables' stored constraints (TR Definition 4.1); run the solver to
a new solved form. 3. \emph{(Close interiors, upward.)} Bind every
remaining variable to its top (\(1_\emptyset\) / the empty row);
re-solve once more (re-emissions against tops discharge by DI-top and TR
Proposition 2.6, so this cannot fail).

The downward leaf choice is local: with fully ground caps it is the join
of the stored lower bounds, i.e.~the \(\sqsubseteq\)-least value
satisfying those caps alone. A row cap with holes is a policy guess,
since filling a hole with \(1_\emptyset\) is generally incomparable with
filling it by a concrete atom; parameter errors are the guard against
making that guess silently.

\textbf{Proposition {[}TR Proposition 6.2{]} (Closing terminates and is
sound).} The interleave terminates, and if no
\(\mathsf{error}\)/\(\mathsf{fail}\) occurs, the resulting total ground
substitution \(\gamma_{\mathrm{close}} \in \mathrm{Sol}(\Phi_0)\).
Ground commitments add no variable-to-variable rank facts, so the
variable--variable fragment keeps its model and bounds all growth (the
deficit potential consults only open-result lineages); a rank-breaking
commitment therefore fails \emph{finitely}, through the absolute
fragment. TR Lemma 5.5 is not needed here.

\textbf{Remark {[}TR Example 6.3{]} (Transitive \(\mathrm{glb}\) and
incomplete leaf closing).} Let \[
\Phi_1 = \{3_b \sqsubseteq \alpha,\ \alpha \sqsubseteq \beta,\ 5_b \sqsubseteq \beta\}
\] with \(\alpha,\beta\) leaves. By TR Definition 4.1,
\(3_b \sqsubseteq \alpha \sqsubseteq \beta\) contributes the lower bound
\(3_b \sqsubseteq \beta\), so
\(\mathrm{glb}(\beta)=3_b \vee 5_b = 1_\emptyset\) by DI-cap. Thus
solving collapses \(\beta\) to \(1_\emptyset\) before leaf closing
commits anything. The solutions of \(\Phi_1\) are exactly
\(\beta = 1_\emptyset\), \(\alpha \in \{3_b, 1_\emptyset\}\), and
deterministic closing selects \(\alpha=3_b,\beta=1_\emptyset\) in every
tested emission order; the branch \(\beta \mapsto 5_b\) is not a legal
close-to-\(\mathrm{glb}\) run.

This order-independence is not completeness. Let \[
\Phi_2 = \{3_b \sqsubseteq \alpha,\ 5_b \sqsubseteq \beta,\ \gamma \sqsubseteq \alpha,\ \gamma \sqsubseteq \beta\}
\] with \(\alpha,\beta\) leaves and \(\gamma\) interior. The store is
satisfiable, e.g.~\(\alpha=\beta=1_\emptyset,\gamma=3_b\). Leaf-downward
closing commits the capped leaves to their saturated GLBs,
\(\alpha \mapsto 3_b\) and \(\beta \mapsto 5_b\); re-solving then pins
\(\gamma\) below two distinct atoms and fails. This is a policy
rejection: satisfying the store requires relaxing at least one capped
leaf upward to \(1_\emptyset\), and the deterministic closing policy
deliberately does not backtrack to larger leaf shapes.

\textbf{Proposition {[}TR Proposition 6.4{]} (Uniform-upward closing
yields the greatest solution).} \(\gamma_\uparrow\) of TR Proposition
5.7 is a solution of \(\Phi_0\) --- unconditionally: the soundness
direction of TR Lemma 5.1 composes through policy steps. Its
\emph{greatestness} needs more care than the original one-line proof
(``TR Proposition 5.7 plus TR Theorem 5.9(1)'') admitted: that route
uses TR Theorem 5.9(1)'s only-if direction, which policy steps break ---
the final configuration encodes the run's placement commitments, so TR
Proposition 5.7(ii) quantifies over a possibly \emph{proper} subset of
\(\mathrm{Sol}(\Phi_0)\). Three correct forms:

\begin{enumerate}
\def\labelenumi{(\roman{enumi})}
\item
  \emph{Policy-relative, marked order.} \(\gamma_\uparrow\) is
  pointwise-greatest among the solutions of the
  \emph{policy-strengthened} system --- \(\Phi_0\) plus the run's
  placement commitments; equivalently, the solutions that extend into
  the final configuration (TR Corollary 5.8's phrase). On
  placement-undiscriminated inputs (conjecturally the whole surface
  fragment) this set is all of \(\mathrm{Sol}(\Phi_0)\) and the
  unqualified claim holds.
\item
  \emph{The unqualified marked claim is false.}
  \(\Phi_0 = \{\langle\rho\rangle \approx [\,]\cdot\diamond\cdot[3,5]\}\):
  the run commits \(\rho \mapsto [\,]\cdot\diamond\cdot[3,5]\)
  (inherit-the-split), so
  \(\gamma_\uparrow(\rho) = [\,]\cdot\diamond\cdot[3,5]\), while
  \(\gamma(\rho) = [3]\cdot\diamond\cdot[5] \in \mathrm{Sol}(\Phi_0)\)
  under TR Definition 3.4's \(\approx\)-semantics is
  \(\sqsubseteq\)-\emph{incomparable} to it (TR Proposition 2.9(ii)).
\item
  \emph{Absolute up to \(\approx\), relative to the rank policy.} For
  every \(\gamma \in \mathrm{Sol}(\Phi_0^{\mathrm{rk}})\) --- \(\Phi_0\)
  plus the run's deferred-equation resolutions, the rank-level policy
  commitments of (c-split)/(d) --- and every \(x \in V\):
  \(\gamma(x) \sqsubseteq^{1} \gamma_\uparrow(x)\) at row sort, plain
  \(\sqsubseteq\) at dimension sort. \textbf{{[}proved --- A.6{]}}
  Placement commitments need no exclusion: \(\sqsubseteq^1\) is
  marker-blind on its lower side, which is exactly why the one-shot
  order is the right comparison. Rank commitments do: a solution
  realizing a foreclosed sporadic/conjugate branch can have smaller rank
  at a variable than \(\gamma_\uparrow\)'s committed value, and no
  marker-blind order repairs a rank gap. On deferral-free runs
  \(\Phi_0^{\mathrm{rk}} = \Phi_0\) and the statement is unconditional.
\end{enumerate}

\textbf{Remark {[}TR Example 3.6; §6{]} (Non-principality witnessed).}
\(\Phi = \{3_b \sqsubseteq \alpha\}\), \(\alpha\) a leaf:
\(\gamma_{\mathrm{close}}(\alpha) = 3_b\);
\(\gamma_\uparrow(\alpha) = 1_\emptyset\); both are solutions, neither
factors through the other by substitution (TR Example 3.6). Closing
selects by policy, justified by allocation (leaves) and parsimony
(interiors), not by entailment.

\subsection{7. Projection inference}\label{projection-inference}

Throughout, fix one operation; its shapes are solved and \textbf{closed}
(ground). Its constraint set \(\Phi_{\mathrm{op}}\) is re-derived with
\textbf{freshened projection identifiers}: each axis of each
participating tensor occurrence carries an id \(p \in P\), injectively,
used by this operation only; \(\mathrm{sz}(p) \in \mathbb{N}_{\ge 1}\)
is the axis's solved size. (The spec's fresh row variables are
eliminated by the local re-solve --- TR Lemma 7.7 --- and its label
variables resolve to operand axes, inheriting ids.)

\textbf{Definition {[}TR Definition 7.1{]} (Projection language).} Atoms
\(q ::= \mathsf{Proj}(p) \mid \mathsf{Sol}(\mathit{idx})\), where
\(\mathit{idx}\) is an externally supplied index: \(\mathsf{Fix}(c)\) or
an external symbol. Equations
\(e ::= \mathsf{Eq}(q_1, q_2) \mid \mathsf{Iter}(q)\), with
\(\mathsf{Eq}\) symmetric.

\textbf{Definition {[}TR Definition 7.2{]} (Elaboration
\(\llbracket\cdot\rrbracket\)).} From the ground \(\Phi_{\mathrm{op}}\):

\begin{itemize}
\tightlist
\item
  P1.
  \(d_1 = d_2 \rightsquigarrow \mathsf{Eq}(\llbracket d_1\rrbracket, \llbracket d_2 \rrbracket)\).
\item
  P2. \(d_{\mathrm{res}} \sqsubseteq d_{\mathrm{op}}\) with
  \(\mathrm{sz}(d_{\mathrm{op}}) = 1 \rightsquigarrow \mathsf{Eq}(\llbracket d_{\mathrm{op}}\rrbracket, \mathsf{Sol}(\mathsf{Fix}\,0))\)
  (sever and pin).
\item
  P3. \(d_{\mathrm{res}} \sqsubseteq d_{\mathrm{op}}\) otherwise
  \(\rightsquigarrow \mathsf{Eq}(\llbracket d_{\mathrm{res}}\rrbracket, \llbracket d_{\mathrm{op}}\rrbracket)\).
\item
  P4. \(R_1 \approx R_2\): outer-edge alignment; P1 per aligned pair
  (ranks match --- shapes are closed and the equality held).
\item
  P5. \(R_{\mathrm{res}} \sqsubseteq R_{\mathrm{op}}\): outer-edge
  alignment; P2/P3 per aligned pair, with P2's severance applied
  row-wise (result side of the pair \(\rightsquigarrow \mathsf{Iter}\),
  operand pinned); each surplus interior result axis
  \(\rightsquigarrow \mathsf{Iter}(\llbracket\cdot\rrbracket)\).
\item
  P6. each terminal axis
  \(\rightsquigarrow \mathsf{Iter}(\llbracket\cdot\rrbracket)\). Side
  condition: \(\mathsf{Sol}(\mathsf{Fix}\,c)\) is emitted only with
  \(0 \le c <\) the axis's size (slices are validated against
  solved sizes).
\end{itemize}

\textbf{Definition {[}TR Definition 7.3{]} (Solver).} A single pass over
the finite equation set \(E\) maintaining: a union--find partition of
\(P\); a partial pin map on classes; an iterate set \(I\) of classes. -
\(\mathsf{Eq}(\mathsf{Proj}\,p, \mathsf{Proj}\,q)\): union,
\(\mathsf{fail}\) if class sizes differ. -
\(\mathsf{Eq}(\mathsf{Proj}\,p, \mathsf{Sol}\,i)\): pin \(p\)'s class at
\(i\), \(\mathsf{fail}\) on a conflicting pin. -
\(\mathsf{Eq}(\mathsf{Sol}\,i, \mathsf{Sol}\,j)\): check \(i = j\).
\(\mathsf{Iter}(\mathsf{Proj}\,p)\): add class to \(I\). -
\(\mathsf{Iter}(\mathsf{Sol}\,\_)\): no-op.

\textbf{Labeling:} pinned class \(\mapsto\) its pin (pins dominate
\(I\)); unpinned class in \(I\) with size \(> 1\) \(\mapsto\) a fresh
iterator symbol; otherwise \(\mapsto \mathsf{Fix}(0)\). The
\textbf{product space} \(P^{\times} = \prod_j R_j\), one factor
\(R_j = \{0..n_j{-}1\}\) per fresh iterator with \(n_j\) the class size;
the \textbf{index map} \(\pi_T\) of tensor \(T\) maps each axis to its
class's label.

\textbf{Theorem {[}TR Theorem 7.4{]} (Canonicity).} For a finite \(E\):
the partition, the pin map, the iterate set, and the labeling are
independent of processing order; the solver fails iff (a) some
\(\approx\)-class (the least equivalence containing the
\(\mathsf{Eq}(\mathsf{Proj},\mathsf{Proj})\) pairs) contains two ids of
different sizes, or (b) some class receives two distinct pins, or (c)
some \(\mathsf{Eq}(\mathsf{Sol}\,i, \mathsf{Sol}\,j)\) has \(i \ne j\).
On success the output is unique up to a bijective renaming of the fresh
iterator symbols.

\textbf{Definition {[}TR Definition 7.5{]} (Reduction; coverage).} A
factor \(j\) of \(P^\times\) is a \emph{reduction axis} iff no component
of \(\pi_{\mathrm{LHS}}\) is \(\mathsf{Iter}(s_j)\).

\textbf{Proposition {[}TR Proposition 7.6{]} (Injectivity, surjectivity,
and the reduction characterization).} With core labels only
(\(\mathsf{Iter}/\mathsf{Fix}\)):

\begin{itemize}
\tightlist
\item
  \begin{enumerate}
  \def\labelenumi{(\roman{enumi})}
  \tightlist
  \item
    \(\pi_{\mathrm{LHS}} : P^\times \to \mathrm{Addr}\) is injective iff
    every product variable occurs in some component of
    \(\pi_{\mathrm{LHS}}\) --- i.e.~iff there is no reduction axis.
  \end{enumerate}
\item
  \begin{enumerate}
  \def\labelenumi{(\roman{enumi})}
  \setcounter{enumi}{1}
  \tightlist
  \item
    \(\pi_{\mathrm{LHS}}\) is surjective onto the LHS's index domain iff
    every LHS axis of size \(> 1\) is labeled \(\mathsf{Iter}\) and the
    map (axis \(\mapsto\) its variable) is injective on those axes.
  \end{enumerate}
\item
  \begin{enumerate}
  \def\labelenumi{(\roman{enumi})}
  \setcounter{enumi}{2}
  \tightlist
  \item
    Hence: the accumulating read-modify-write is required iff a
    reduction axis exists (non-injectivity), and pre-initialization is
    required iff \(\pi_{\mathrm{LHS}}\) is non-surjective or (under an
    erasing accumulator) non-injective --- the
    \passthrough{\lstinline!=:+!} analysis, now a corollary.
  \end{enumerate}
\end{itemize}

\textbf{Lemma {[}TR Lemma 7.7{]} (Local re-solve succeeds; locality).}
If the global solve and closing succeeded, then for each operation: (a)
the per-operation re-derivation (spec template with fresh row/dimension
variables against the ground operand shapes) has a solution, and the
core solver finds it, eliminating every spec variable; (b) the resulting
\(\approx\) depends only on \(\Phi_{\mathrm{op}}\) --- by construction
of the freshened ids, no constraint of any other operation mentions any
id of \(P\), so no external fact can enter the closure.

\textbf{Theorem {[}TR Theorem 7.8{]} (Soundness of projections w.r.t.
shapes).} Let shapes be solved and closed, the per-operation elaboration
as in TR Definition 7.2 with its side condition, and the local solve
succeed. Then: 1. \emph{(Bounded addressing)} for every participating
tensor \(T\) and every \(p \in P^\times\), \(\pi_T(p)\) is a valid
multi-index: componentwise, an \(\mathsf{Iter}(s_j)\) component
addresses an axis whose size equals \(n_j\), and \(s_j < n_j\); a
\(\mathsf{Fix}(c)\) component has \(c <\) the axis size. 2.
\emph{(Co-iteration is finer than size-equality)}
\(p \approx q \Rightarrow \mathrm{sz}(p) = \mathrm{sz}(q)\), and
\(\approx\) is the \textbf{least} equivalence justified by this
operation's own constraints (TR Theorem 7.4 + TR Lemma 7.7(b)) --- in
particular, axes identified in size by other operations' constraints are
not co-iterated here.

\end{document}
